\documentclass[uplatex,dvipdfmx,ja=standard,a4paper,11pt]{bxjsarticle}
\usepackage{geometry}
\geometry{margin=25mm}

\begin{document}

\begin{center}
    {\huge 生成AI利用申告書}
\end{center}

\vspace{1.5em}

\noindent
25G1110 早川 和輝 \\
2026年4月28日 \\
使用した AIツール：Gemini

\vspace{2em}

\noindent
今回の課題において，ArduinoとProcessingを連携させてマイクの音信号の波形を描画させるプログラムを作成した際，プログラムが正しく連動して動作するかの特定の確認に生成AIを使用した．

\vspace{1em}

\noindent
作成の手順として，まず自身で作成したArduinoとProcessingのソースコード全体をプロンプトに入力した上で，課題の要件とともに「このプログラムは正しく動きますか」という質問と指示を出した．本件は「一部修正して使った」という区分に該当する．これは，AIから提示された修正案のうち，シリアル通信を正しく行うためのボーレート（通信速度）の設定値を一致させるという指摘を採用し，自身のプログラムの数値を書き換えたためである．

\vspace{1em}

\noindent
具体的には，AIに指摘されたシリアル通信のボーレートの不一致の箇所を修正し，Processing側の設定値を115200からArduino側と同じ921600に変更してコードを修正した後に実際にプログラムを実行した．その後，意図した通りにマイク入力からの波形がProcessingのウィンドウに描画されることを目視で確認した．

\vspace{1em}

\noindent
AIの回答に関する間違いはなかったと判断した．AIが指摘した通り，シリアル通信においては送信側と受信側でボーレートが合致していないとデータが正しく読み取れないという仕様は客観的な事実であり，実際にボーレートの数値を揃えた結果として波形が正常に表示されたためである．

\vspace{1em}

\noindent
今回の利用において，AIからは通信エラーの原因であるボーレートの不一致の指摘と具体的な修正箇所のみが提案された．私は通信の不具合を解消することを最優先事項とし，AIが提案した通信速度の修正方式を自身の判断で選択して実装した．

\vfill
\begin{center}
1
\end{center}

\end{document}
